\subsection*{Endpoint reduction for the two true A.11 correction supports}
After the angular interpolation the unedited source azimuthal profile
$E_b(y)$, $y=\log X$, is independent of $\eta$. On the interval
$[y_-,y_+]$ between the extreme support endpoints of the two
width-$3/10$ A.11 correction bumps, its logarithmic slope is
$\ell_b=\partial_y\log(\sqrt X E_b)=-\lambda$. Write the actual
corrected profile as
\[
 E_c(y,\eta)=E_b(y)(1+u(y,\eta)),\qquad
 u=\sum_{i=1}^2c_i(\eta)\beta(y-y_i).
\]
The two source bumps have disjoint supports and differ from the
axial moment corrections (A.9).

\begin{proposition}[Full-support corrected-field bounds]
\label{prop:a11-endpoint-jets}
Suppose the genuine moment-correction coefficients obey the source
bounds $|c_i|\le192\lambda^{28}$ and
$|c_{i,\eta}|\le12\lambda^{28}$ and the prescribed smooth-step
derivative satisfies $|\beta|\le640/3$, $|\beta'|\le200000$.
For $0<\lambda\le1/20$ put
$\delta=40960\lambda^{28}<10^{-31}$. Then, simultaneously at
every point of both supports and for every $|\eta|\le1$,
\begin{align*}
 (1-\delta)E_b&\le E_c\le(1+\delta)E_b,\\
 |\partial_\eta\log E_c|&\le
     2560\lambda^{28}/(1-\delta),\\
 |\ell_c+\lambda|&\le38400000\lambda^{28}/(1-\delta)
     <\lambda/2,\\
 |\partial_\eta\ell_c|&\le
     2400000\lambda^{28}/(1-\delta)^2.
\end{align*}
\end{proposition}
\begin{proof}
At most one bump is active at a given radius.
For its coefficient $c$ and bump $\beta$, the exact formulas are
\[
 \partial_\eta\log E_c=\frac{c_\eta\beta}{1+c\beta},
 \quad \ell_c=-\lambda+\frac{c\beta'}{1+c\beta},
 \quad\partial_\eta\ell_c=\frac{c_\eta\beta'}{(1+c\beta)^2}.
\]
The last identity includes an exact numerator cancellation.
Insert the specified source bounds and
$40960/20^{28}<10^{-31}$.
\end{proof}

\begin{theorem}[Positive-tail comparison and endpoint extrema]
\label{thm:a11-endpoint-tails}
Assume the actual terminal $S(\infty)=0$ moment condition and
pressure gauge $\Pi(\infty)=0$, and define
\[
 D_k=\frac1{2XE_k^2}\int_X^\infty E_k(x)^2dx,\qquad
 P_k=\frac1{2E_k^2}\int_X^\infty\frac{E_k(x)^2}{x}dx,
 \quad k=b,c.
\]
For any $y$ on either correction support, with
$\kappa=(1+\delta)/(1-\delta)$,
\[
 \kappa^{-2}D_b(y)\le D_c(y,\eta)\le\kappa^2D_b(y),
 \quad
 \kappa^{-2}P_b(y)\le P_c(y,\eta)\le\kappa^2P_b(y).
\]
Here $\kappa^2\le1+8\delta$. Moreover
\[
 D_b'=2\lambda D_b-\tfrac12,\qquad
 P_b'=(1+2\lambda)P_b-\tfrac12
\]
on all of $[y_-,y_+]$; hence the maximum of each uncorrected
tail moment on both supports occurs at an extreme endpoint.
Writing $D_*=\max(D_b(y_-),D_b(y_+))$ and
$P_*=\max(P_b(y_-),P_b(y_+))$, one obtains
\[
 E_c\le(1+\delta)E_b(y_-),\quad
 D_c\le(1+8\delta)D_*,\quad
 P_c\le(1+8\delta)P_*.
\]
If the \emph{entire corrected future} has $\ell_c\le0$, its
positive pressure kernel gives the independent $P_c\le1/2$.
\end{theorem}
\begin{proof}
The quotient of corrected and uncorrected normalized future
tail kernels is exactly
$\bigl[(1+u(y+t,\eta))/(1+u(y,\eta))\bigr]^2$.
It lies in $[\kappa^{-2},\kappa^2]$ because
$|u|\le\delta$ at every future radius. Positivity of the
integrals gives the comparison. Differentiating the unedited
tails gives $D_b'=-1/2-2\ell_bD_b$ and
$P_b'=-1/2+(1-2\ell_b)P_b$. With $\ell_b=-\lambda$,
their solutions are monotone or constant. Also $E_b$ decreases
strictly on $[y_-,y_+]$ and
$\kappa^2=1+4\delta/(1-\delta)^2\le1+8\delta$.
\end{proof}

\begin{corollary}[A three-value conditional source gate]
\label{cor:a11-three-values}
Let $\varepsilon=2560\lambda^{28}/(1-\delta)$.
If the same actual admissible source tuple and true terminal
normalizations additionally establish
\[
\frac{(1+\delta)E_b(y_-)}{\lambda}\le\frac14,\quad
4h(1+8\delta)D_*\le\frac{61}{100},\quad
2\varepsilon((1+8\delta)D_*+\tfrac12)\le\frac1{100},\quad
2A\le\frac{1001}{1000},
\]
and all corrected future slopes are nonpositive, then
$|N_s|/E_c\le1621\lambda/4000<\lambda/2$ with
exact slack $379\lambda/4000$ on both supports. The further
stress-cone statement needs separate $Q_s$ and radius bounds.
\end{corollary}

\paragraph{Source authentication remains open.}
This reduces the supportwise task to $E_b(y_-)$,
$D_b(y_-),D_b(y_+)$ and the true coefficient and
terminal conditions. It does not supply certified endpoint
values for a single actual Appendix-A tuple, nor prove the
outer stress cone, Appendix-B contraction, distinct unforced
startup solutions or physical zero-force residual. Finite
rational or synthetic interval tests do not fill these gaps.
